
\documentclass[../../ps_q.tex]{subfiles}
    
\begin{document}

\section{運動方程式と束縛条件}

\subsection{地表付近での重力}

図のように，鉛直面内において水平右向きに$x$軸，鉛直上向きに$y$軸を定める．
原点Oから，初速$v_0$，$x$軸の正方向からの仰角$\theta$で，小物体（質量$m$）を投射した．
重力加速度の大きさを$g$とする．

小物体の加速度の$x$成分$a_x$，$y$成分$a_y$を求めよ．

\vspace{4pt}
\begin{center}

    \begin{tikzpicture}
        \draw[->,>=stealth,semithick] (0,0) -- (7,0) node[right] {$x$};
        \draw[->,>=stealth,semithick] (0,0) -- (0,5) node[above] {$y$};
        \draw[->,>-stealth,thick] (0,0) -- (1,1.5) node[above right=-2pt] {$v_0$};
        \draw (1,0) coordinate (A) -- (0,0) coordinate (O) -- (0.2,0.3) coordinate (B);
        \filldraw[thick,fill=black!50!white] (0,0) circle (5pt) node[below left] {O};
        \draw pic["$\theta$", draw=black, angle eccentricity=1.35, angle radius=7mm] 
            {angle=A--O--B};

    \end{tikzpicture}
\end{center}

\end{document}
